\documentclass[11pt,a4paper]{article}

\usepackage[margin=1in]{geometry}
\usepackage{amsmath,amssymb,mathtools}
\usepackage{enumitem}

\newcommand{\StartE}{\mathit{StartE}}
\newcommand{\EndE}{\mathit{EndE}}
\newcommand{\NextB}{\mathit{Next}_{B}}
\newcommand{\GType}{\mathit{GType}}
\newcommand{\gtype}{\mathit{gtype}}

\title{Core BPMN Process-State Semantics}
\author{}
\date{}

\begin{document}
\maketitle

\section{BPMN Process Structure}

A BPMN process is a tuple
\[
\boxed{B=(T,F)}
\]
where
\[
\boxed{T=E\uplus A\uplus G}
\]
is a finite set of process transitions, consisting of Events $E$, Activities $A$, and Gateways $G$. Events are partitioned into Start Events and End Events:
\[
\boxed{E=\StartE\uplus\EndE.}
\]
The sequence-flow relation is
\[
\boxed{F\subseteq T\times T.}
\]

\section{Predecessors and Successors}

For each transition $t\in T$, its predecessor and successor sets are defined by
\[
\boxed{{}^{\bullet}t=\{t'\in T\mid (t',t)\in F\}}
\]
and
\[
\boxed{t^{\bullet}=\{t'\in T\mid (t,t')\in F\}.}
\]

\section{Structural Well-Formedness}

The BPMN profile satisfies the following structural constraints.

\paragraph{Start Events.}
\[
\boxed{
\forall s\in\StartE:\quad
{}^{\bullet}s=\varnothing
\;\land\;
|s^{\bullet}|=1
}
\tag{WF-S}
\]

\paragraph{End Events.}
\[
\boxed{
\forall e\in\EndE:\quad
|{}^{\bullet}e|\geq 1
\;\land\;
e^{\bullet}=\varnothing
}
\tag{WF-E}
\]

\paragraph{Activities.}
\[
\boxed{
\forall a\in A:\quad
|{}^{\bullet}a|\geq 1
\;\land\;
|a^{\bullet}|=1
}
\tag{WF-A}
\]

\paragraph{Gateways.}
\[
\boxed{
\forall g\in G:\quad
|{}^{\bullet}g|\geq 1
\;\land\;
|g^{\bullet}|\geq 1
}
\tag{WF-G1}
\]
A gateway with exactly one incoming and one outgoing sequence flow is excluded:
\[
\boxed{
\forall g\in G:\quad
|{}^{\bullet}g|>1
\;\lor\;
|g^{\bullet}|>1
}
\tag{WF-G2}
\]

\section{Gateway Types}

The gateway types considered in the current profile are
\[
\boxed{
\GType=\{\mathit{Exclusive},\mathit{Inclusive},\mathit{Parallel}\}.
}
\]
Each gateway has exactly one type:
\[
\boxed{
\gtype:G\rightarrow\GType.
}
\]
The operational effect of each gateway type is captured later by $\NextB$.

\section{BPMN Process State}

A transition configuration is a non-empty set
\[
\boxed{\varnothing\neq\sigma_T\subseteq T.}
\]
A BPMN process state is denoted by
\[
\boxed{m(\sigma_T).}
\]
For ordinary transitions,
\[
\boxed{
t\in\sigma_T
\iff
t\text{ is currently activated and ready to be executed at }m(\sigma_T).
}
\]
An End Event contained in $\sigma_T$ denotes that the corresponding process branch has reached an End Event; End Events are terminal and are not executed further.

\section{Initial and Final States}

All Start Events constitute the initial transition configuration:
\[
\boxed{\sigma_T^{0}=\StartE}
\]
and therefore
\[
\boxed{m_0=m(\StartE).}
\]
No empty initial marking is required.

A state is considered a final process state exactly when all End Events have been reached:
\[
\boxed{
\mathit{Final}_B\bigl(m(\sigma_T)\bigr)
\iff
\EndE\subseteq\sigma_T.
}
\tag{FINAL}
\]
Before all End Events have been reached, an individual End Event may coexist with other activated transitions. Once the process is final, no non-End transition may remain activated:
\[
\boxed{
\EndE\subseteq\sigma_T
\;\Rightarrow\;
\sigma_T=\EndE.
}
\tag{WF-FINAL}
\]

\section{General State-Transition Rule}

Let
\[
\boxed{
\NextB(t,\sigma_T)\subseteq t^{\bullet}
}
\]
denote the set of successor transitions activated as a consequence of executing $t$. Its precise definition is left to the operational semantics of Activities and Gateway types.

The general process-transition rule is
\[
\boxed{
\frac{t\in\sigma_T\setminus\EndE}
{m(\sigma_T)
\xrightarrow{\;t\;}_{B}
 m\!\left((\sigma_T\setminus\{t\})\cup\NextB(t,\sigma_T)\right)}
}
\tag{EXEC}
\]
Thus, executing $t$ removes $t$ from the current configuration, preserves all other currently activated transitions, and activates the transitions returned by $\NextB(t,\sigma_T)$.

\section{Process State Space}

The BPMN process-state space is
\[
\boxed{
\Sigma_B
=
\{m(\sigma_T)\mid \varnothing\neq\sigma_T\subseteq T\}.
}
\]
Only states satisfying the well-formedness conditions above are considered valid process states.

\end{document}
